% from ega3-1-fr.tex
\oldpage[0\textsubscript{III}]{55}
Avec les notations de (\hyperref[0.12.1.3-fr]{12.1.3}), supposons que l'on
ait un diagramme commutatif
\[
\label{0.12.1.6.1-fr}
\tag{12.1.6.1}
\xymatrix{
0\ar[r]&\mathcal F\ar[r]^r\ar[d]_u
&\mathcal G\ar[r]^s\ar[d]_v
&\mathcal H\ar[r]\ar[d]_w&0\\
0\ar[r]&\mathcal F'\ar[r]_{r'}
&\mathcal G'\ar[r]_{s'}
&\mathcal H'\ar[r]&0
}
\]
où $r$, $s$ sont des homomorphismes de $\mathcal O_X$-Modules, $r'$, $s'$ des
homomorphismes de $\mathcal O_{X'}$-Modules, $u$, $v$, $w$ des $f$-morphismes
et les lignes sont \emph{exactes}. On en déduit alors un diagramme commutatif
\[
\label{0.12.1.6.2-fr}
\tag{12.1.6.2}
\xymatrix@C=6pt{
\cdots\ar[r]&H^p(X,\mathcal F)\ar[r]\ar[d]_{u_p}
&H^p(X,\mathcal G)\ar[r]\ar[d]_{v_p}
&H^p(X,\mathcal H)\ar[r]^\partial\ar[d]_{w_p}
&H^{p+1}(X,\mathcal F)\ar[r]\ar[d]_{u_{p+1}}&\cdots\\
\cdots\ar[r]&H^p(X',\mathcal F')\ar[r]
&H^p(X',\mathcal G')\ar[r]
&H^p(X',\mathcal H')\ar[r]_\partial
&H^{p+1}(X',\mathcal F')\ar[r]&\cdots
}
\]
En effet, (\hyperref[0.12.1.6.1-fr]{12.1.6.1}) se factorise en
\[
\xymatrix@R=14pt{
0\ar[r]&\mathcal F\ar[r]\ar[d]
&\mathcal G\ar[r]\ar[d]
&\mathcal H\ar[r]\ar[d]&0\\
0\ar[r]&\psi^*(\mathcal F)\ar[r]\ar[d]
&\psi^*(\mathcal G)\ar[r]\ar[d]
&\psi^*(\mathcal H)\ar[r]\ar[d]&0\\
0\ar[r]&\mathcal F'\ar[r]
&\mathcal G'\ar[r]
&\mathcal H'\ar[r]&0
}
\]
où la ligne du milieu est \emph{exacte}
(\hyperref[0.3.7.2-fr]{0\textsubscript{I}, 3.7.2}) et il suffit d'utiliser le
fait que (\hyperref[0.12.1.3.2-fr]{12.1.3.2}) est un homomorphisme de
$\delta$-foncteurs et que les $H^p(X',\mathcal F')$ forment un
$\delta$-foncteur en $\mathcal F'$.
\end{remark}

\begin{env}[12.1.7]
\label{0.12.1.7-fr}
Les hypothèses et notations étant celles de
(\hyperref[0.12.1.3-fr]{12.1.3}), considérons maintenant le cas où
$\mathcal F=f_*(\mathcal F')=\psi_*(\mathcal F')$; nous allons voir que le
di-homomorphisme défini dans (\hyperref[0.12.1.3-fr]{12.1.3})
\[
\label{0.12.1.7.1-fr}
\tag{12.1.7.1}
H^p(X,f_*(\mathcal F'))\longrightarrow H^p(X',\mathcal F')
\]
peut s'obtenir (à un automorphisme près de $H^p(X',\mathcal F')$) comme
\emph{edge-homomorphisme d'une suite spectrale} du foncteur composé
$\mathcal F'\leadsto\Gamma(X,\psi_*(\mathcal F'))$ (T, 2.4). La description
de l'homomorphisme (\hyperref[0.12.1.7.1-fr]{12.1.7.1}) donnée dans
(\hyperref[0.12.1.4-fr]{12.1.4}) montre ici qu'on peut obtenir cet
homomorphisme de la façon suivante : on considère des résolutions injectives
$\mathcal L^\bullet$ et $\mathcal L'^\bullet$ de $\psi_*(\mathcal F')$ et de
$\mathcal F'$ respectivement, puis on prend un homomorphisme de complexes
$v:\psi^*(\mathcal L^\bullet)\to\mathcal L'^\bullet$ «~au-dessus~» de
l'homomorphisme canonique $\psi^*(\psi_*(\mathcal F'))\to\mathcal F'$;
